\chapter{直线方程与线性方程组}\label{ch:g10:lines}

坐标平面上的直线可以用一个\omterm{def:g10:algebra:equation}{方程}描述，两条直线
合在一起构成一个\omterm{def:g10:algebra:equation}{方程}组，其解即为它们的
交点。本章联系三种语言——几何（直线）
、代数（\omterm{def:g10:algebra:equation}{方程}）和\omterm{def:g10:functions:function}{函数}（\cref{ch:g10:reffunc}中的\omterm{def:g10:reffunc:affine}{一次函数}）
——并介绍求解\omterm{def:g10:lines:system}{线性方程组}的
两种标准方法。

\section{直线的方程}

\begin{theorem}[斜截式方程]\label{thm:g10:lines:reduced}
在\omterm{def:g10:coordgeom:system}{坐标系}中：
\begin{enumerate}
  \item 每条非竖直的直线都有形如
        \[
        y = mx + p ,
        \]
        其中 $m$（\emph{斜率}\index{斜率}）与 $p$（
        \emph{\omterm{def:g10:reffunc:affine}{纵截距}}，即 $y$ 轴上的截距）被唯一确定——这就是该直线的
        \emph{斜截式方程}\index{斜截式方程}；
  \item 每条竖直的直线都有\omterm{def:g10:algebra:equation}{方程} $x = c$；
  \item 一点位于直线上，当且仅当其坐标满足
        该\omterm{def:g10:algebra:equation}{方程}。
\end{enumerate}
\end{theorem}
\admitted

\begin{proposition}[过两点的斜率]\label{prop:g10:lines:slope}
过 $A(x_A, y_A)$ 与 $B(x_B, y_B)$ 的（非竖直）直线，
其中 $x_A \neq x_B$，其\omterm{def:g10:reffunc:affine}{斜率}为
\[
m = \frac{y_B - y_A}{x_B - x_A}
\qquad
\text{（“$y$ 的变化量与 $x$ 的变化量之比”）。}
\]
\end{proposition}

\begin{proof}
两点都满足 $y = mx + p$，故 $y_B - y_A = (mx_B + p) - (mx_A + p)
= m(x_B - x_A)$；两边除以 $x_B - x_A \neq 0$。
\end{proof}

\begin{omfigure}
\begin{tikzpicture}
\begin{axis}[omaxis,
  width=8.2cm, height=5.8cm,
  xmin=-1.4, xmax=4.4, ymin=-1.4, ymax=4.4,
  xtick={1,2,3}, ytick={1,3}]
  \addplot[omDef, thick, domain=-1.1:4.1] {x + 0.5};
  \addplot[omThm, thick, domain=-1.1:4.1] {3};
  \draw[omMeth, thick] (axis cs:2,-1.2) -- (axis cs:2,4.2);
  \node[omDef, font=\small, rotate=38] at (axis cs:2.7,3.6)
    {$y = x + 0.5$};
  \node[omThm, font=\small] at (axis cs:-0.55,3.3) {$y = 3$};
  \node[omMeth, font=\small, right] at (axis cs:2,-0.8) {$x = 2$};
\end{axis}
\end{tikzpicture}

{\small 三类直线：斜直线（$m \neq 0$）、水平直线（$m = 0$），
竖直直线（没有\omterm{thm:g10:lines:reduced}{斜截式方程}——\omterm{def:g10:algebra:equation}{方程}为 $x = c$）。}
\end{omfigure}

\begin{method}[求过两点的直线方程]\label{met:g10:lines:through}
已知 $A$ 与 $B$，其中 $x_A \neq x_B$：
\begin{enumerate}
  \item 计算\omterm{def:g10:reffunc:affine}{斜率} $m = \dfrac{y_B - y_A}{x_B - x_A}$；
  \item 写出 $y = mx + p$，代入 $A$（或 $B$）的坐标，
        求出 $p$；
  \item 用\emph{另一个}点检验所得的\omterm{def:g10:algebra:equation}{方程}。
\end{enumerate}
若 $x_A = x_B$，则该直线为竖直直线，其\omterm{def:g10:algebra:equation}{方程}为 $x = x_A$。
\end{method}

\begin{example}\label{ex:g10:lines:through}
过 $A(1, 5)$ 与 $B(3, 1)$ 的直线。\omterm{def:g10:reffunc:affine}{斜率}：
$m = \frac{1 - 5}{3 - 1} = \frac{-4}{2} = -2$。于是 $y = -2x + p$；由
点 $A$ 得 $5 = -2 \times 1 + p$，故 $p = 7$。所求\omterm{def:g10:algebra:equation}{方程}为
$y = -2x + 7$。代入 $B$ 检验：$-2 \times 3 + 7 = 1$。
\end{example}

\begin{proposition}[平行直线]\label{prop:g10:lines:parallel}
两条非竖直的直线平行，当且仅当它们
\omterm{def:g10:reffunc:affine}{斜率}相同。两条不同的平行直线没有公共点；\omterm{def:g10:reffunc:affine}{斜率}不同的两条直线
恰有一个公共点。
\end{proposition}

\begin{proof}
设 $y = mx + p$ 与 $y = m'x + p'$ 为两条直线。公共点
满足 $mx + p = m'x + p'$，即 $(m - m')x = p' - p$。若 $m \neq m'$，
则该\omterm{def:g10:algebra:equation}{方程}恰有一个解 $x = \frac{p' - p}{m - m'}$，即一个
交点。若 $m = m'$，\omterm{def:g10:algebra:equation}{方程}化为 $0 = p' - p$：
当 $p \neq p'$ 时无解（两条不同的平行线）；对每个 $x$ 都成立的情形为 $p = p'$
（为同一条直线）。因此，平行（即没有交点，或两直线重合）
等价于 $m = m'$。
\end{proof}

\begin{remark}[方向向量]
直线 $y = mx + p$ 每向右移动一个单位，就向上移动 $m$ 个单位，因此
\omterm{def:g10:vectors:vector}{向量} $\vec u\,(1, m)$ 是该直线的一个\emph{方向
向量}\index{方向向量}，即该直线与 $\vec u$ 平行。
两条直线平行，当且仅当它们的方向\omterm{def:g10:vectors:vector}{向量}
\omterm{def:g10:vectors:collinear}{共线}（见\cref{thm:g10:vectors:det}）。
\end{remark}

\section{线性方程组}

\begin{definition}[线性方程组]\label{def:g10:lines:system}
关于未知数 $x$ 与 $y$ 的一个\emph{二元一次方程组}\index{线性方程组}，
是指如下的一对\omterm{def:g10:algebra:equation}{方程}
\[
\begin{cases}
a x + b y = c \\
a' x + b' y = c'
\end{cases}
\]
两个\omterm{def:g10:algebra:equation}{方程}须\emph{同时}成立。\emph{解}是指满足两个\omterm{def:g10:algebra:equation}{方程}的一对
$(x, y)$。每个\omterm{def:g10:algebra:equation}{方程}都描述一条直线，因此求解该
\omterm{def:g10:algebra:equation}{方程}组就是求两条直线的交点。
\end{definition}

\begin{method}[代入法]\label{met:g10:lines:substitution}
\begin{enumerate}
  \item 从一个\omterm{def:g10:algebra:equation}{方程}中解出一个未知数——选最简单的，例如
        将 $y$ 用 $x$ 表示；
  \item 将该式代入\emph{另一个}\omterm{def:g10:algebra:equation}{方程}，
        此时\omterm{def:g10:algebra:equation}{方程}中只含 $x$；
  \item 解出 $x$，再由第 1 步求出 $y$；
  \item 将所得的一对数代入两个原\omterm{def:g10:algebra:equation}{方程}进行检验。
\end{enumerate}
\end{method}

\begin{example}\label{ex:g10:lines:substitution}
用代入法解\omterm{def:g10:algebra:equation}{方程}组 $\begin{cases} x + 2y = 8 \\ 3x - y = 3 \end{cases}$。
由第二个\omterm{def:g10:algebra:equation}{方程}得 $y = 3x - 3$。将其代入
第一个\omterm{def:g10:algebra:equation}{方程}：
\[
x + 2(3x - 3) = 8
\ \Longleftrightarrow\
x + 6x - 6 = 8
\ \Longleftrightarrow\
7x = 14
\ \Longleftrightarrow\
x = 2 ,
\]
则 $y = 3 \times 2 - 3 = 3$，所以\omterm{def:g10:algebra:equation}{方程}组的解为数对 $(2, 3)$。
检验：$2 + 2 \times 3 = 8$，$3 \times 2 - 3 = 3$。
\end{example}

\begin{method}[消元法]\label{met:g10:lines:elimination}
\begin{enumerate}
  \item 将每个\omterm{def:g10:algebra:equation}{方程}乘以适当的数，使某个未知数
        在两个\omterm{def:g10:algebra:equation}{方程}中的系数互为相反数；
  \item 将两个\omterm{def:g10:algebra:equation}{方程}相加，消去该未知数；
  \item 解出剩下的一元一次\omterm{def:g10:algebra:equation}{方程}，再代回
        求出另一个未知数；
  \item 将所得的一对数代入两个原\omterm{def:g10:algebra:equation}{方程}进行检验。
\end{enumerate}
\end{method}

\begin{example}\label{ex:g10:lines:elimination}
用消元法解\omterm{def:g10:algebra:equation}{方程}组 $\begin{cases} 2x + 3y = 7 \\ 5x - 2y = 8 \end{cases}$。
为消去 $y$，将第一个\omterm{def:g10:algebra:equation}{方程}乘以 $2$，第二个
\omterm{def:g10:algebra:equation}{方程}乘以 $3$：
\[
\begin{cases}
4x + 6y = 14 \\
15x - 6y = 24
\end{cases}
\]
两式相加得 $19x = 38$，故 $x = 2$。将其代入 $2x + 3y = 7$，
得 $4 + 3y = 7$，故 $y = 1$。解为 $(2, 1)$。验算第二个
\omterm{def:g10:algebra:equation}{方程}：$5 \times 2 - 2 \times 1 = 8$。
\end{example}

\begin{omfigure}
\begin{tikzpicture}
\begin{axis}[omaxis,
  width=8.2cm, height=6.0cm,
  xmin=-0.9, xmax=4.9, ymin=-1.4, ymax=4.9,
  xtick={2}, ytick={1}]
  \addplot[omDef, thick, domain=-0.6:4.6] {(7 - 2*x)/3};
  \addplot[omThm, thick, domain=-0.5:2.9] {(5*x - 8)/2};
  \fill (axis cs:2,1) circle (2pt);
  \node[above right, font=\small] at (axis cs:2,1) {$(2, 1)$};
  \node[omDef, font=\small] at (axis cs:4.05,0.35) {$2x + 3y = 7$};
  \node[omThm, font=\small] at (axis cs:1.3,3.3) {$5x - 2y = 8$};
\end{axis}
\end{tikzpicture}

{\small 从几何角度看\omterm{def:g10:lines:system}{线性方程组}：每个\omterm{def:g10:algebra:equation}{方程}都是一条直线，而
\cref{ex:g10:lines:elimination}中的解 $(2,1)$ 是它们的
交点。}
\end{omfigure}

\begin{remark}[有几个解？]
与两条直线类似，一般地，\omterm{def:g10:lines:system}{线性方程组}恰有一个解
（两直线不平行），特殊情形下无解（两条不同的平行线）或
有无穷多解（两个\omterm{def:g10:algebra:equation}{方程}为同一条直线）。判别准则见
\cref{thm:g10:vectors:det}，\omterm{def:g10:algebra:equation}{方程}组有唯一解当且仅当
$ab' - a'b \neq 0$。
\end{remark}

\begin{example}[列方程组建模]\label{ex:g10:lines:modeling}
三本笔记本和两支钢笔共 $8$ 元；一本笔记本和四支钢笔共
$6$ 元。设 $x$ 为一本笔记本的价格，$y$ 为一支钢笔的价格：
\[
\begin{cases}
3x + 2y = 8 \\
x + 4y = 6 .
\end{cases}
\]
由第二个\omterm{def:g10:algebra:equation}{方程}得 $x = 6 - 4y$；代入得
$3(6 - 4y) + 2y = 8$，即 $18 - 10y = 8$，得 $y = 1$，进而
得 $x = 2$。一本笔记本 $2$ 元，一支钢笔 $1$ 元。
\end{example}

\section{练习}

\begin{exercise}[$\star$]\label{exo:g10:lines:1}
对每条直线，读出其\omterm{def:g10:reffunc:affine}{斜率}和 $y$ 轴截距，并画出草图：
\[
y = 2x - 3, \qquad
y = -\tfrac13 x + 1, \qquad
y = 4, \qquad
x = -2 .
\]
\end{exercise}

\begin{exercise}[$\star$]\label{exo:g10:lines:2}
判断点 $A(2, 7)$ 是否在直线 $y = 3x + 1$ 上？点 $B(-1, -1)$ 呢？
点 $C(4, 13)$ 呢？
\end{exercise}

\begin{exercise}[$\star$]\label{exo:g10:lines:3}
求过下列两点的直线的\omterm{thm:g10:lines:reduced}{斜截式方程}：
\[
\text{(a) } A(0, 2) \text{ 和 } B(3, 8);
\qquad
\text{(b) } A(-1, 5) \text{ 和 } B(2, -4);
\qquad
\text{(c) } A(4, 1) \text{ 和 } B(4, 6).
\]
\end{exercise}

\begin{exercise}[$\star$]\label{exo:g10:lines:4}
在直线 $y = 3x - 1$、$y = -3x + 2$、$y = 3x + 5$、
$y = \frac{6x + 2}{2}$ 中，哪些直线互相平行？哪些
实际上是同一条直线？
\end{exercise}

\begin{exercise}[$\star$]\label{exo:g10:lines:5}
用代入法解下列\omterm{def:g10:algebra:equation}{方程}组：
\[
\begin{cases}
y = 2x - 1 \\
3x + y = 9
\end{cases}
\qquad\text{然后}\qquad
\begin{cases}
x - y = 4 \\
2x + 3y = 3 .
\end{cases}
\]
\end{exercise}

\begin{exercise}[$\star$]\label{exo:g10:lines:6}
用加减消元法解下列\omterm{def:g10:algebra:equation}{方程}组：
\[
\begin{cases}
x + y = 10 \\
x - y = 4
\end{cases}
\qquad\text{然后}\qquad
\begin{cases}
3x + 2y = 12 \\
2x + 5y = 8 .
\end{cases}
\]
\end{exercise}

\begin{exercise}[$\star\star$]\label{exo:g10:lines:7}
不解\omterm{def:g10:algebra:equation}{方程}组，判断下列各\omterm{def:g10:algebra:equation}{方程}组各有多少组解：
\[
\begin{cases}
2x - y = 3 \\
-4x + 2y = -6
\end{cases}
\qquad
\begin{cases}
2x - y = 3 \\
-4x + 2y = 1
\end{cases}
\qquad
\begin{cases}
2x - y = 3 \\
x + y = 0 .
\end{cases}
\]
\end{exercise}

\begin{exercise}[$\star\star$]\label{exo:g10:lines:8}
一场校园演出共售出两百张票，其中一部分票价为 $5$（儿童票），
另一部分票价为 $9$（成人票），总额为 $1352$。每种票各
售出多少张？
\end{exercise}

\begin{exercise}[$\star\star$]\label{exo:g10:lines:9}
设直线 $d$ 为 $y = 2x - 3$，点 $A(1, 4)$。
\begin{enumerate}
  \item 求直线 $d'$ 的\omterm{def:g10:algebra:equation}{方程}，它平行于 $d$，
        且经过 $A$。
  \item 求 $d'$ 与 $x$ 轴的交点。
\end{enumerate}
\end{exercise}

\begin{exercise}[$\star\star\star$]\label{exo:g10:lines:10}
设 $A(0, 0)$、$B(6, 2)$ 和 $C(2, 4)$。
\begin{enumerate}
  \item 求三角形 $ABC$ 中由顶点 $A$ 和 $B$ 引出的中线的\omterm{def:g10:algebra:equation}{方程}，
        即中线连接一个顶点与其对边的
        \omterm{prop:g10:coordgeom:midpoint}{中点}。
  \item 求它们的交点 $G$，并验证
        $G$ 也在第三条中线上。（你应当发现
        $G$ 的坐标是 $A$、$B$、$C$ 坐标的平均数。）
\end{enumerate}
\end{exercise}

\section{问题：唯一解、无解还是无穷多解}

\begin{problem}[{周末问题——两条直线有三种可能的命运，
行列式决定其归属，而世界上最古老的教科书
早已知道此事}]
\label{pb:g10:lines:1}
\omterm{def:g10:lines:system}{二元一次方程组}，两个未知数：它们所画出的一对直线
可能交于一点、永不相交，或本为同一条直线——而
\emph{每一个}\omterm{def:g10:lines:system}{线性方程组}都继承了这三种命运之一。
本题对它们分类，并引入决定性的数（
\omterm{def:g10:vectors:vector}{向量}一章中的旧相识），并解答
两千年前中国古代的雉兔同笼问题，它见于
\emph{《九章算术》}，并以区域的角点作结——这是
今天每家航空公司和工厂所用优化方法的第一
步。

\medskip
\noindent\textbf{第一部分——直线，信手拈来。}
\begin{enumerate}
  \item 求过
        $(1, 2)$ 与 $(3, 8)$ 的直线的\omterm{thm:g10:lines:reduced}{斜截式方程}
        （见\cref{met:g10:lines:through}）。
  \item 在 $y = 3x - 1$、$y = 3x + 4$ 与 $y = -x + 7$ 中，哪两条直线
        互相平行，
        （见\cref{prop:g10:lines:parallel}），其余两对直线在何处相交？
        （求出其中一个交点。）
  \item 过 $(2, 1)$ 与 $(2, 5)$ 的直线，为何没有
        形如 $y = mx + p$ 的\omterm{thm:g10:lines:reduced}{斜截式方程}？它的
        \omterm{def:g10:algebra:equation}{方程}是什么？
  \item 点 $(4, 11)$ 是否在直线 $y = 3x - 1$ 上？
        点 $(-1, -5)$ 呢？
  \item 过点 $(0, 3)$ 且\omterm{def:g10:reffunc:affine}{斜率}为 $-2$ 的直线，求它的
        \omterm{def:g10:algebra:equation}{方程}、它在两条坐标轴上的截距，以及
        它在第一象限内截出的三角形的面积。
\end{enumerate}

\medskip
\noindent\textbf{第二部分——三种命运。}
\begin{enumerate}[resume]
  \item 用代入法求解
        （见\cref{met:g10:lines:substitution}）：
        $\begin{cases} y = 2x - 3 \\ 3x + y = 7 \end{cases}$
  \item 用加减消元法求解
        （见\cref{met:g10:lines:elimination}）：
        $\begin{cases} 2x + 3y = 8 \\ 5x - 3y = -1 \end{cases}$
  \item 对下列两个\omterm{def:g10:algebra:equation}{方程}组进行分类，
        并借助直线加以解释：
        \[
        \begin{cases} 2x - y = 3 \\ 4x - 2y = 6 \end{cases}
        \qquad\text{ 和 }\qquad
        \begin{cases} 2x - y = 3 \\ 4x - 2y = 1 . \end{cases}
        \]
        写出完整的三分情形：相交、平行、
        重合——一解、无解、无穷多解。
  \item 对一般\omterm{def:g10:algebra:equation}{方程}组 $ax + by = e$，$cx + dy = f$，
        决定性的数是 $ad - bc$——与
        \omterm{def:g10:vectors:collinear}{共线}性判别所用的行列式相同（见\cref{pb:g10:vectors:1}）。
        请解释这一巧合（哪两个\omterm{def:g10:vectors:vector}{向量}
        \omterm{def:g10:vectors:collinear}{共线}，恰好当且仅当两条直线
        平行或重合？），并计算
        第 7 题与第 8 题中三个\omterm{def:g10:algebra:equation}{方程}组的 $ad - bc$。
  \item 当 $m$ 取何值时，
        $\begin{cases} x + my = 3 \\ 2x + 6y = 7 \end{cases}$
        该\omterm{def:g10:algebra:equation}{方程}组无解？请对其余所有
        $m$\,\dots\ 的值求解该\omterm{def:g10:algebra:equation}{方程}组，或至少说明你会如何求解。
\end{enumerate}

\medskip
\noindent\textbf{第三部分——《九章算术》。}
\begin{enumerate}[resume]
  \item 取自\emph{《九章算术》}
        （中国，约两千年前）：一只笼子里装着雉和兔，
        共 $35$ 个头、$94$ 条腿。
        问各有多少只？
  \item 一个果汁摊将果汁含量为 $60\,\%$ 的糖浆与
        果汁含量为 $20\,\%$ 的饮料混合，得到 $30$ 升、果汁含量为 $35\,\%$ 的饮品，
        两种原料各需多少升？
  \item 一个两位数的各位数字之和为 $12$；交换其十位与个位数字后，
        所得之数比原数减少 $18$。求此数，
        并利用初中分册中“数字代数”周末问题的方法列出\omterm{def:g10:algebra:equation}{方程}组。
  \item 面包房里，$3$ 杯咖啡与 $2$ 个羊角面包共需
        $9.10$ 欧元；$2$ 杯咖啡与 $3$ 个羊角面包共需
        $8.90$ 欧元。请巧妙求解：将两个\omterm{def:g10:algebra:equation}{方程}\emph{相加}
        与\emph{相减}，各能
        直接告诉你什么？
  \item 一个集市摊位声称：$2$ 个苹果 $+ 1$ 个梨售价
        $5$ 欧元，$4$ 个苹果 $+ 2$ 个梨售价 $9$ 欧元。
        求解——或者更确切地说，请说明该\omterm{def:g10:algebra:equation}{方程}组的
        “命运”揭示了该摊位的算术有什么问题。
\end{enumerate}

\medskip
\noindent\textbf{第四部分——角点说了算。}
\begin{enumerate}[resume]
  \item 直线 $y = 3x - 1$ 将平面分为两部分。原点位于
        哪一侧？试描述集合
        $y > 3x - 1$ 所表示的区域，并说明如何检验一个点是否满足该不等式。
  \item 画出由下列四个约束条件所确定的区域
        \[
        x \geq 0, \qquad y \geq 0, \qquad x + y \leq 4,
        \qquad y \leq 2x + 1,
        \]
        并求出它的四个角点。
  \item 某车间的利润为 $P = 3x + 2y$，定义在
        第 17 题的区域上。计算 $P$ 在四个角点处的值，并在本阶段不加证明地承认
        此类线性量在此类区域上的\omterm{def:g10:functions:extrema}{最大值}总能在角点处取得
        （想象直线
        $3x + 2y =$ 常数 扫过区域），给出最优
        生产方案。
  \item 将初中分册“双温度计”周末问题中的两辆出租车
        写成\omterm{def:g10:algebra:equation}{方程}组并求解：
        $\begin{cases} y = 1.5x + 3 \\ y = 2x \end{cases}$，并
        解释其解的意义。
  \item 终章，即“词典”：一个\omterm{def:g10:algebra:equation}{方程} $\leftrightarrow$
        一条直线；一个\omterm{def:g10:algebra:equation}{方程}组 $\leftrightarrow$ 两条直线，
        三种命运，由 $ad - bc$ 决定；多个不等式
        $\leftrightarrow$ 一个区域，其角点承载
        最优值。用三句话概括——你刚刚
        参观了\emph{线性代数}与
        \emph{线性规划}的门厅，两者都将在
        大学分册中完整建立。
\end{enumerate}
\end{problem}
